%============================================================= §6 OBSERVED
\section{Observed Output Across Machines \& Empirical Verification Matrix}

Requirement~(c) asks what was observed on each machine involved in the reconnaissance and attack
pipeline. Because our toolkit operates as a command-line reconnaissance and exploitation engine
whose primary outputs are structured terminal telemetry, the verbatim console records below
represent the decisive artifacts of system behavior across both offensive and defensive roles.

\subsection{Attacker Workstation Observations}
On the attacking machine (\texttt{192.168.0.105}), the console provides real-time visibility into
each attack phase:
\begin{enumerate}
    \item \textbf{Phase 1 (SYN Half-Open Scan):} Displays instantaneous port classification
    (\textsc{open}, \textsc{closed}, \textsc{filtered}) alongside the extracted initial TTL and
    TCP Window Size.
    \item \textbf{Phase 2 (Banner Grabbing):} Prints extracted service strings and daemon version
    metadata.
    \item \textbf{Phase 3 (OS Fingerprinting):} Displays the probabilistic verdict, confidence tier,
    and a bulleted breakdown of supporting evidence.
    \item \textbf{Phase 4 \& 5 (Weaponization \& Exploitation):} Reports payload dispatch status,
    exfiltrated file previews, and remote code execution outputs.
    \item \textbf{Severance Observation (Defense Run):} When active defense is enabled on the victim,
    the attacker console observes connection timeouts (\texttt{WinError 10060}) mid-execution once
    firewall quarantine rules take effect.
\end{enumerate}

\subsection{Victim Server Daemon Logs}
The victim machine hosts the vulnerable lab services (\texttt{victim/http\_server.py} on port~80 and
\texttt{victim/ftp\_server.py} on port~21). The service logs below record the corresponding inbound
exploit interactions from the server's perspective:

\begin{logbox}
[Victim Terminal 1: HTTP Server (Port 80)]
[HTTP] Vulnerable HTTP server listening on 0.0.0.0:80 (CGI enabled for uploads/*.py)
[HTTP] ('192.168.0.105', 51482) -> HEAD /
[HTTP] ('192.168.0.105', 51488) -> GET /uploads/shell.py?cmd=whoami
[HTTP] ('192.168.0.105', 51488) -> CGI executed shell.py (query='cmd=whoami')
[HTTP] ('192.168.0.105', 51493) -> GET /..\..\..\..\..\windows\win.ini  (200 OK, served)
[HTTP] ('192.168.0.105', 51498) -> GET /..%5c..%5c..%5cwindows%5cwin.ini (200 OK, served)

[Victim Terminal 2: FTP Server (Port 21)]
[FTP] Vulnerable FTP server listening on 0.0.0.0:21
[FTP] 220 LabFTP Server Ready
[FTP] USER anonymous -> 331 Password required
[FTP] PASS anon@lab -> 230 Login successful
[FTP] PASV -> 227 Entering Passive Mode (192,168,0,116,195,84)
[FTP] LIST -> 150 Here comes the directory listing
[FTP] STOR uploads/shell.py -> 226 Transfer complete
[FTP] DELE uploads/shell.py -> 250 File deleted successfully (Forensic cleanup)
\end{logbox}

\subsection{Master Summary Matrix of Empirical Verification Runs}
Table~\ref{tab:master_matrix} summarizes the definitive empirical verification runs archived in
\texttt{final-log/}, cross-referencing target platforms, tested services, and observed defense states.

\begin{table}[htbp]
\centering
\footnotesize
\renewcommand{\arraystretch}{1.18}
\caption{Consolidated Master Matrix of Empirical Verification Test Runs}
\label{tab:master_matrix}
\begin{tabularx}{\textwidth}{@{}L{2.8cm} L{2.6cm} L{3.0cm} L{2.2cm} X@{}}
\toprule
\hdrrow \thh{Log Name} & \thh{Target System} & \thh{Tested Ports \& States} & \thh{Defense State} & \thh{Attribution \& Outcome}\\
\midrule
\texttt{full-chain.log} & Windows 11 Desktop & 21, 80 (Open), 22, 81 (Filtered) & Active Blocker On & Windows 62\%; 3/4 exploits succeeded; severed on Run 2.\\
\addlinespace
\texttt{cross-platform.log} & Windows 11 + Kali 6.x & 21, 80 (Win) / 21, 22, 80 (Kali) & Passive / Off & Dynamic OS payload swap: \texttt{win.ini} vs.\ \texttt{/etc/passwd}.\\
\addlinespace
\texttt{kali.log} & Kali Linux VM & 22, 80, 81 (All Closed / RST) & Passive / Off & Linux 76\% (Closed-port RST) / 67\% (ICMP Code 9 fallback).\\
\addlinespace
\texttt{scan\_100\_ports.log} & Windows 10/11 Target & 100 ports (2 $\times$ 50 batches) & Passive / Off & Batch 1: Win 71\% (SYN/ACK); Batch 2: Win 69\% (RST).\\
\addlinespace
\texttt{ipad.log} & Apple iPad (iPadOS 17) & 80 (Open VLC), 8080 (Closed) & Passive / Off & macOS/BSD 48\% (Darwin option order, WScale=5).\\
\addlinespace
\texttt{phn.log} & Android 14 Smartphone & 2221 (Open FTP), 2122 (Closed) & Passive / Off & Linux 58\% (SACK before TS, RFC 6864 IP ID=0).\\
\addlinespace
\texttt{ssh\_attack.log} & Windows OpenSSH Host & 22 (Open SSH) & Passive / Off & Windows 75\%; cracked \texttt{test:test123}; loud desktop payloads.\\
\addlinespace
\texttt{defense\_win.log} & Windows 11 Desktop & Multi-port scan (Scenarios A--G) & Windows Firewall & Tripped at 8 SYNs; canary port 81; auto-release after 60s.\\
\addlinespace
\texttt{defense\_kali.log} & Kali Linux VM & Multi-port scan (Scenarios A--G) & Linux Netfilter & Netfilter \texttt{iptables} DROP counter climbing dynamically.\\
\bottomrule
\end{tabularx}
\end{table}

\subsection{Cross-Platform Operating System Comparison}
Table~\ref{tab:os_comparison} highlights the contrasting stack signatures observed between Microsoft
Windows and Unix-like targets across identical probe sequences:

\begin{table}[htbp]
\centering
\small
\renewcommand{\arraystretch}{1.2}
\caption{Cross-Platform Network Stack Behavioral Signatures Observed in Testbed}
\label{tab:os_comparison}
\begin{tabularx}{\textwidth}{@{}L{3.6cm} L{4.6cm} X@{}}
\toprule
\hdrrow \thh{Diagnostic Metric} & \thh{Windows 11 Desktop} & \thh{Kali Linux / Android}\\
\midrule
Initial IPv4 TTL & 128 & 64\\
TCP Window Size & 65535 or 8192 & 29200, 65160, or 0 (closed RST)\\
IP-ID Sequence Archetype & \texttt{STRICT\_INCREMENTAL} ($\Delta \approx 1\dots 3$) & \texttt{ALL\_ZERO} ($\mathrm{ID} = 0$, RFC 6864 DF=1)\\
TCP Timestamp Clock Rate & Calibrated $\approx 1000$\,Hz (1\,ms tick) & $\approx 250$\,Hz or 1000\,Hz\\
RFC 3168 ECN Negotiation & Ignored (\texttt{0x12} returned) & Supported (\texttt{0x52} ECE reflected)\\
Invalid Flag Probe (SYN+FIN) & Emits TCP RST on open port & Silently dropped without reply\\
ICMP Non-Standard Code 9 & Normalized to Code 0 & Reflected unaltered as Code 9\\
OS-Aware Traversal Target & \texttt{windows\textbackslash win.ini} & \texttt{/etc/passwd}\\
RCE Command Execution & \texttt{whoami} $\to$ \texttt{homed01\textbackslash turjo} & \texttt{id} $\to$ \texttt{uid=0(root)}\\
\bottomrule
\end{tabularx}
\end{table}
